\documentclass[11pt]{article}
\usepackage{fontspec}
\usepackage[french,english]{babel}
\usepackage{amsmath}
\usepackage{amssymb}
\usepackage{amsthm}
\usepackage{mathrsfs}
\usepackage{tikz-cd}
\usepackage{enumitem}

\newcommand{\printedpage}[1]{\marginpar{\footnotesize p.~#1}}
\newcommand{\msp}[1]{\marginpar{\footnotesize [#1]}}
\newcommand{\uHom}{\underline{\mathrm{Hom}}}
\newcommand{\Ens}{\mathrm{Ens}}
\newcommand{\Fais}{\mathrm{Fais}}
\newcommand{\pstop}{\mathrm{pstop}}

\begin{document}

\printedpage{} % no page number printed
% head: (none)

\begin{center}
{\Large\textsf{THE DERIVATORS}}

\bigskip

\textbf{ALEXANDRE GROTHENDIECK}

\bigskip\bigskip

\textbf{Chapter XVIII}

\medskip

\textbf{Accessible categories and sets}
\end{center}

\newpage

\printedpage{} % no page number printed
% head: (none)

This text was deciphered and transcribed in \LaTeXe{} by M. Künzer. It was edited by M. Künzer, J. Malgoire, and G. Maltsiniotis. The transcription is as faithful as possible to the manuscript. Only a few obvious minor corrections have been made. For the rare additions or comments by the editors, as well as for the original numbering of the manuscript pages, typewriter characters \texttt{[typewriter]} in brackets are used. A few footnotes added by the editors are indicated by ``~Ed.~N.~''. The page numbering of this chapter begins at page 50, because originally this chapter was the continuation of chapter XIX which ends in the middle of page 50.

This edition is provisional. Remarks, comments and corrections are welcome. Send a message to:

\textbf{[e-mail address removed]}

\hfill G. Maltsiniotis\hspace*{2cm}

\newpage

\printedpage{1}
% head: A. Grothendieck : Dérivateurs, Ch. XVIII, Édition des Universités de Montpellier II et Paris VII
% foot: Version 17 décembre 2025

\msp{page 50}

\section*{4\quad Paratopos and the operation $\mathcal{M}\boxtimes\mathcal{N}$}

\subsection*{4.1\quad Pseudo-topos. Problem of the product.}

I call a \textit{pseudo-$\mathfrak{U}$-topos category}, or simply a \textit{pseudo-topos}, a $\mathfrak{U}$-category $\mathcal{M}$ stable under $\mathfrak{U}$-direct limits, and which admits a \textit{small} subcategory $\mathcal{M}_0$ that is \textit{strictly generating} (i.e. generating by strict epimorphisms). Then $\mathcal{M}$ also admits $\mathfrak{U}$-inverse limits (SGA 4 I 8.12.7). (But the example of (Gr), cf. loc. cit. 8.12.9, shows that $\mathcal{M}^\circ$ is not necessarily, for all that, a pseudo-$\mathfrak{U}$-topos, for lack of $\mathcal{M}$ admitting a small subcategory that is strictly \textit{co}generating. But theorem 2.6 \texttt{[of chapter XIX]} gives conditions on $\mathcal{M}$ (mainly of the exactness type) % [sic] par forcément -> pas forcément, misprint; also the parenthesis opened at "(Mais" is never closed
\msp{page 51}
for $\mathcal{M}$ to admit enough injectives -- which will then imply the existence of a small strictly cogenerating subcategory, hence that $\mathcal{M}^\circ$ is also a pseudo-topos (\footnote{This is the case in particular if $\mathcal{M}$ is a topos.}).

I visualize pseudo-topos categories as being the category of ``sheaves'' on an object $\mathcal{X}$, visualized as a sort of ``space'' which I distinguish from the category $\mathcal{M}=\Fais(\mathcal{X})$ of sheaves on $\mathcal{X}$. It is $\mathcal{X}$ that will be for this the ``pseudo-topos''. Pseudo-topoi form a 2-category, the variance in $\mathcal{X}$ being opposite to that in $\Fais(\mathcal{X})=\mathcal{M}$. To be precise, if $\mathcal{X}$ corresponds to $\mathcal{M}$, $\mathcal{Y}$ to $\mathcal{N}$, a morphism $f$ from $\mathcal{X}$ to $\mathcal{Y}$ will be given by a functor
\[
f^* \;:\; \mathcal{N} \longrightarrow \mathcal{M}
\tag*{(4.1.1)}
\]
in the opposite direction, commuting with arbitrary direct $\varinjlim$ (\footnote{Say here that the objects of $\mathcal{M}$ are visualized as generalized ``opens'' of $\mathcal{X}$. We have $\mathcal{M}\xrightarrow{\sim}\uHom^!(\mathcal{M}^\circ,\Ens)$, i.e. $\mathcal{M}$ is indeed interpreted as the ``sheaves'' on this category $\mathcal{M}$ of opens \ldots}). (\textbf{N.B.} if $\mathcal{X}$, $\mathcal{Y}$ are even topoi, i.e. $\mathcal{M}$, $\mathcal{N}$ topos categories, and if we want $f$ to be a morphism of topoi, we will require in addition that $f^*$ be in addition \textit{left exact}. We will say that a \textit{morphism of pseudo-topoi is exact} if in addition $f^*$ is left exact.)

\msp{page 52}
By virtue of the dual of loc. cit. 8.12.7, this is equivalent to the fact that $f^*$ admits a right adjoint $f_*$,
\[
f_* \;:\; \mathcal{M} \longrightarrow \mathcal{N}\;,
\tag*{(4.1.2)}
\]
which therefore commutes with small $\varprojlim$\,. The datum of $f_*$ satisfying this property therefore determines $f^*$ up to unique isomorphism. But in general, it is not enough to give oneself such an $f_*$, to be able to assert that it corresponds to an $f^*$, i.e. that it admits a right adjoint\,; for this, we would have to know that the functors $\mathcal{M}\longrightarrow\Ens$ which commute with small $\varprojlim$ are representable, which will be the case if $\mathcal{M}$ admits a small subcategory that is \textit{co}generating, i.e. if $\mathcal{M}^\circ$ is also a pseudo-topos category. % [sic] qu'il admette un adjoint à droite -> adjoint à gauche, f^* is the left adjoint of f_*

\newpage

\printedpage{2}
% head: A. Grothendieck : Dérivateurs, Ch. XVIII, Édition des Universités de Montpellier II et Paris VII
% foot: Version 17 décembre 2025

It is probably wise to proceed as in the case of topoi, and to call a \textit{morphism of pseudo-topoi} $f:\mathcal{X}\longrightarrow\mathcal{Y}$ a quadruple
\[
(f^*,f_*,\alpha,\beta)
\tag*{(4.1.3)}
\]
\msp{page 53}
of two adjoint functors $(f^*,f_*)$, and of the adjunction morphisms. When $f^*$ is an equivalence, which is equivalent to $f_*$ an equivalence, or $\alpha$ and $\beta$ isomorphisms, we will say that we have an \textit{equivalence of pseudo-topoi}, and for all practical purposes, we will then be able to regard $\mathcal{X}$ and $\mathcal{Y}$ as essentially identical objects.

For two pseudo-topoi $\mathcal{X}$, $\mathcal{Y}$, I feel like introducing (as in the case of topoi) a product pseudo-topos
\[
\mathcal{X}\times\mathcal{Y}\;,
\tag*{(4.1.4)}
\]
which would have the properties of a 2-product in the 2-category of pseudo-topoi (\footnote{It turns out in the end that what we construct is \textit{not} a 2-product in the 2-category under consideration, so it is better to denote it by a sign such as $\mathcal{X}\boxtimes\mathcal{Y}$. But if $\mathcal{X}$, $\mathcal{Y}$ are two topoi, $\mathcal{X}\boxtimes\mathcal{Y}$ is a 2-product in the 2-category of topoi.}). I should specify that by definition
\[
\begin{aligned}
\uHom_{\pstop}(\mathcal{X},\mathcal{Y}) \;&\simeq\; \uHom_!(\mathcal{F}(\mathcal{Y}),\mathcal{F}(\mathcal{X}))^\circ\\
&\simeq\; \uHom^{!'}(\mathcal{F}(\mathcal{X}),\mathcal{F}(\mathcal{Y})),
\end{aligned}
\tag*{(4.1.5)}
\]
where the subscript ! (resp. the superscript !) indicates the full subcategory in the category $\uHom$ under consideration consisting of the functors which commute with small $\varinjlim$ (resp. with small
\msp{page 54}
$\varprojlim$\,), and the accent in the $\uHom^{!'}$ in the last member indicates that we require in addition that the functor in question have a left adjoint.

So let
\[
\mathcal{M} \;=\; \mathcal{F}(\mathcal{X})\;, \qquad \mathcal{N} \;=\; \mathcal{F}(\mathcal{Y})\;,
\tag*{(4.1.6)}
\]
and let $\mathcal{Z}$ be a third pseudo-topos, described by
\[
\mathcal{P} \;=\; \mathcal{F}(\mathcal{Z})\;.
\tag*{(4.1.7)}
\]
We want to make explicit, up to equivalence, the product category $\uHom_{\pstop}(\mathcal{Z},\mathcal{X})\times\uHom_{\pstop}(\mathcal{Z},\mathcal{Y})$, i.e.
\[
\uHom_!(\mathcal{M},\mathcal{P})^\circ \;\times\; \uHom_!(\mathcal{N},\mathcal{Q})^\circ.
\tag*{(4.1.8)}
\] % [sic] \mathcal{Q} -> \mathcal{P}, \mathcal{Q} is not defined; \mathcal{Z} is described by \mathcal{P} (4.1.7)
I am going to introduce for this a category
\printedpage{3}
% head: A. Grothendieck : Dérivateurs, Ch. XVIII, Édition des Universités de Montpellier II et Paris VII
% foot: Version 17 décembre 2025
\msp{page 55}
\[
\mathcal{M}\boxtimes\mathcal{N}
\tag*{(4.1.9)}
\]
by the formula
\[
\mathcal{M}\boxtimes\mathcal{N} \;\simeq\; \uHom^{!!}(\mathcal{M}^\circ\times\mathcal{N}^\circ,\Ens)\;,
\tag*{(4.1.10)}
\]
where the sign $\uHom^{!!}$ denotes the full subcategory of the category $\uHom$, consisting of the bi-contrafunctors % doubt resolved: retain the diplomatic bi-contrafoncteurs reading; see corrections.csv, PDF leaf 5
\[
F \;:\; \mathcal{M}^\circ\times\mathcal{N}^\circ \longrightarrow \Ens\;,
\tag*{(4.1.10.1)}
\]
such that for every $x$ in $\mathcal{M}$, the functor
\[
y \longmapsto F(x,y) \;:\; \mathcal{N}^\circ \longrightarrow \Ens
\tag*{(4.1.10.2)}
\]
is in $\uHom^!(\mathcal{N}^\circ,\Ens)$, i.e. transforms small $\varinjlim$ into $\varprojlim$\,, and that, symmetrically, for every $y$ in $\mathcal{N}$, the functor
\[
x \longmapsto F(x,y) \;:\; \mathcal{M}^\circ \longrightarrow \Ens
\tag*{(4.1.10.3)}
\]
is in $\uHom^!(\mathcal{M}^\circ,\Ens)$. But this also means that these functors are \textit{representable}\,; so we also find
\[
\mathcal{M}\boxtimes\mathcal{N} \xleftarrow{\sim} \uHom^!(\mathcal{N}^\circ,\mathcal{M})
\tag*{(4.1.11)}
\]
and
\[
\mathcal{M}\boxtimes\mathcal{N} \xleftarrow{\sim} \uHom^!(\mathcal{M}^\circ,\mathcal{N})
\tag*{(4.1.12)}
\]
\msp{page 56}
using (in each of these formulas) on both sides the equivalences of categories
\[
\left\{
\begin{aligned}
&\mathcal{M} &&\xrightarrow{\sim} &&\uHom^!(\mathcal{M}^\circ,\Ens)\\
&\mathcal{N} &&\xrightarrow{\sim} &&\uHom^!(\mathcal{N}^\circ,\Ens)\;.
\end{aligned}
\right.
\tag*{(4.1.13)}
\]
We can also say that the category $\mathcal{M}\boxtimes\mathcal{N}$ is equivalent to that of \textit{pairs} of functors
\[
\left\{
\begin{aligned}
&\varphi &&: &&\mathcal{M}^\circ &&\longrightarrow &&\mathcal{N}\\
&\psi &&: &&\mathcal{N}^\circ &&\longrightarrow &&\mathcal{M}\;,
\end{aligned}
\right.
\tag*{(4.1.14)}
\]
equipped with two adjunction morphisms
\[
\left\{
\begin{aligned}
&\psi\varphi^\circ &&\xleftarrow{\alpha} &&\mathrm{id}_{\mathcal{M}}\\
&\varphi\psi^\circ &&\xleftarrow{\beta} &&\mathrm{id}_{\mathcal{N}}
\end{aligned}
\right.
\tag*{(4.1.15)}
\]
satisfying the two usual compatibilities expressing the adjunction of $\varphi^\circ$ and of $\psi$ (or, equivalently, of $\psi^\circ$ and of $\varphi$). The functor
\[
(\varphi,\psi) \longmapsto \varphi \qquad (\text{resp.}\;\; (\varphi,\psi) \longmapsto \psi)
\]
gives an equivalence of the category of these pairs with $\uHom^!(\mathcal{M}^\circ,\mathcal{N})$ (resp. with $\uHom^!(\mathcal{N}^\circ,\mathcal{M})$), so that we have a commutative square of canonical equivalences of functors:

\end{document}
